\chapter{Détermination de structure~: RMN du \texorpdfstring{$^{13}$C}{13C}, SM et IR réunies}\label{ch:b2:structure-determination}

Un laboratoire de parfumerie reçoit un flacon de liquide incolore qui sent le jasmin. Personne n'y a
collé d'étiquette. Un \omterm{def:b2:mass-spec-atomic:spectrum}{spectre de masse} donne sa formule brute, un spectre infrarouge ses groupes
fonctionnels, un spectre du carbone~13 son squelette et un spectre du proton la façon dont les fragments
sont reliés~; en une heure, le liquide a un nom. Ce chapitre ajoute la \omterm{def:b2:structure-determination:c13}{RMN du carbone~13} à la RMN du
proton et à la spectroscopie infrarouge du volume de première année et à la \omterm{def:b2:mass-spec-atomic:spectrum}{spectrométrie de masse} du
\cref{ch:b2:mass-spec-atomic}, et fait des quatre une seule méthode.

\begin{recall}
Le volume de première année~: déplacement chimique et blindage, protons équivalents, intégration,
couplage spin--spin et règle des $n + 1$ uplets, nombres d'onde infrarouges des groupes fonctionnels et
empreinte digitale, nombre d'insaturations, résolution d'une structure à partir de trois spectres.
\Cref{ch:b2:mass-spec-atomic}~: \omterm{def:b2:mass-spec-atomic:molecular-ion}{ion moléculaire}, \omterm{def:b2:mass-spec-atomic:isotope-pattern}{massifs isotopiques}, fragmentation.
\end{recall}

\section{La RMN du carbone~13}

Le carbone~12, isotope principal, n'a pas de spin nucléaire et ne donne aucun signal de RMN. Le
carbone~13 a un spin
$\tfrac12$, comme le proton, mais il ne représente que 1.1~\% du carbone naturel et son signal est plus
faible, noyau pour noyau, que celui d'un proton~; un spectre du carbone exige donc davantage
d'échantillon ou davantage d'acquisitions.
% ledger: iso:C13

\begin{definition}[RMN du carbone~13]\label{def:b2:structure-determination:c13}
La \emph{RMN du carbone~13}\index{RMN du carbone 13} enregistre les résonances des noyaux \ce{^{13}C}
d'un échantillon. On l'effectue le plus souvent avec \emph{découplage large bande}\index{decouplage large bande@découplage large bande}~: les protons sont
irradiés sur toute leur gamme de fréquences pendant la mesure, ce qui supprime tous les couplages
carbone--proton, si bien que chaque sorte de carbone donne une seule raie. Des \emph{carbones
équivalents}\index{carbones equivalents@carbones équivalents}, échangés par une symétrie de la molécule
(ou par une rotation rapide), donnent la même raie.
\end{definition}

\begin{proposition}[Pas de couplage carbone--carbone]\label{prop:b2:structure-determination:no-cc-coupling}
Les couplages entre atomes de carbone voisins n'apparaissent pas dans un spectre du \ce{^{13}C} en
abondance naturelle.
\end{proposition}

\begin{proof}
Un couplage entre deux carbones ne s'observe que dans les molécules où tous deux sont des \ce{^{13}C}.
Pour deux positions voisines données, la probabilité vaut $a_{13}^2 = 0.011^2 \approx 1.2 \times 10^{-4}$~:
environ une molécule sur huit mille. Le signal de chaque carbone provient presque entièrement de molécules
dans lesquelles ses voisins sont des \ce{^{12}C}, invisibles pour le spectromètre~; les satellites couplés
sont cent fois plus faibles que le bruit d'un spectre ordinaire.
\end{proof}

\begin{proposition}[Nombre de raies]\label{prop:b2:structure-determination:signals}
Un spectre du \ce{^{13}C} découplé présente une raie par ensemble de \omterm{def:b2:structure-determination:c13}{carbones équivalents}, sauf si deux
ensembles ont par hasard le même déplacement.
\end{proposition}

\begin{proof}[Argument]
Des \omterm{def:b2:structure-determination:c13}{carbones équivalents} se trouvent dans des environnements identiques et ont donc le même
déplacement~; des carbones non équivalents diffèrent en général. Les couplages avec les protons étant
supprimés et les couplages entre carbones absents
(\cref{prop:b2:structure-determination:no-cc-coupling}), rien ne dédouble une raie. Une superposition
accidentelle se produit quand deux carbones différents ont des déplacements plus proches que la
résolution, ce qui est courant pour les carbones d'un cycle benzénique vers \qty{128}{ppm}.
\end{proof}

Les déplacements s'étendent sur environ \qty{220}{ppm}, vingt fois le domaine des protons, si bien que
les superpositions sont plus rares que dans les spectres du proton, et que le déplacement seul indique
la sorte de carbone. Le diagramme ci-dessous montre les déplacements mesurés pour les composés de ce chapitre.

\begin{omfigure}
\begin{tikzpicture}
\begin{axis}[width=0.86\linewidth, height=4.6cm, xmin=0, xmax=220, ymin=-0.6, ymax=4.6, x dir=reverse,
  xlabel={$\delta$ (ppm)}, label style={font=\small}, tick label style={font=\scriptsize},
  ytick={0,1,2,3,4}, yticklabels={C alkyle, C--O, C aromatique, C=O d'ester, C=O de cétone},
  xtick={0,20,40,60,80,100,120,140,160,180,200,220}, grid=major, grid style={black!10}]
\addplot[only marks, mark=|, mark size=5pt, thick, omDef] table[x=shift, y=cls] {figdata/out/bachelor-2-nmr-spectra-d.dat};
\end{axis}
\end{tikzpicture}

{\small Déplacements \ce{^{13}C} mesurés de la butan-2-one, de la pentan-2-one, du benzoate d'éthyle et
de l'éthanoate de benzyle, classés par sorte de carbone (CDCl$_3$~; données PubChem). Les carbonyles de
cétone se situent vers 209, les carbonyles d'ester vers
167--171, les carbones \omterm{def:b2:huckel:aromatic}{aromatiques} entre 128 et 137, les carbones liés à un oxygène d'ester vers 61--66,
et les carbones alkyle au-dessous de 46.}
% figdata: bachelor-2/nmr-spectra.py part d
% ledger: c13:butanone, c13:pentan-2-one, c13:ethylbenzoate, c13:benzylacetate
\end{omfigure}

\begin{proposition}[Hauteur n'est pas nombre]\label{prop:b2:structure-determination:intensities}
Dans un spectre du \ce{^{13}C} découplé de routine, les hauteurs des raies ne sont pas proportionnelles
aux nombres de carbones qu'elles représentent~; les carbones sans hydrogène donnent des raies faibles.
\end{proposition}

\begin{proof}[Argument]
Les spectres de routine répètent la mesure plus vite que les carbones les plus lents ne reviennent à
l'équilibre, et les carbones qui ne portent pas de protons relaxent lentement, si bien que leur signal est
en partie saturé. Le découplage renforce aussi les signaux des carbones porteurs de protons, d'une quantité
qui dépend de leur environnement. Dans la butan-2-one, la raie du carbonyle est la plus faible des quatre,
bien qu'elle représente un seul carbone comme les autres. L'intégration, si utile pour les protons, n'est
donc pas utilisée pour les spectres du carbone de routine.
\end{proof}

\section{La séquence DEPT}

\begin{definition}[DEPT]\label{def:b2:structure-determination:dept}
La \emph{DEPT}\index{DEPT} (\textit{distortionless enhancement by polarisation transfer}, renforcement sans
distorsion par transfert de polarisation) est un ensemble d'expériences du carbone~13 qui donnent aux
signaux des carbones porteurs de protons un signe ou une absence selon leur nombre d'hydrogènes. Un
\emph{carbone quaternaire}\index{carbone quaternaire} ne porte aucun hydrogène (lié à quatre autres
atomes, ou carbone de carbonyle ou \omterm{def:b2:huckel:aromatic}{aromatique} sans H)~; il ne donne aucun signal en DEPT.
\end{definition}

\begin{proposition}[Lire une DEPT]\label{prop:b2:structure-determination:dept}
En DEPT-135, les carbones \ce{CH} et \ce{CH3} pointent vers le haut, les carbones \ce{CH2} vers le bas
et les \omterm{def:b2:structure-determination:dept}{carbones quaternaires} sont absents. En DEPT-90, seuls les carbones \ce{CH} apparaissent.
\end{proposition}

\admitted

L'expérience transfère l'aimantation des protons au carbone auquel ils sont liés, par une séquence
d'impulsions de radiofréquence dont le dernier angle sélectionne la réponse~; la façon dont les
impulsions y parviennent est expliquée dans le volume de troisième année. Comparer le spectre découplé
aux spectres DEPT-135 et DEPT-90 classe chaque raie en C, CH, \ce{CH2} ou \ce{CH3}.

\begin{omfigure}
\pgfplotsset{dept/.style={width=\linewidth, height=2.9cm, ycomb, x dir=reverse, ymin=-1, ymax=1.15,
  ytick=\empty, axis y line=none, axis x line=middle, x axis line style={-}, tick label style={font=\tiny},
  title style={font=\scriptsize, at={(0.0,0.95)}, anchor=north west}, every axis plot/.append style={thick}}}
\begin{minipage}[t]{0.49\linewidth}\centering
{\small butan-2-one}\par
\begin{tikzpicture}
\begin{axis}[dept, xmin=0, xmax=220, xtick={0,50,100,150,200}, title={découplé}]
\addplot[omDef] table[x=shift, y=dec] {figdata/out/bachelor-2-nmr-spectra-a.dat};
\end{axis}
\end{tikzpicture}\par
\begin{tikzpicture}
\begin{axis}[dept, xmin=0, xmax=220, xtick={0,50,100,150,200}, title={DEPT-135}]
\addplot[omThm] table[x=shift, y=d135] {figdata/out/bachelor-2-nmr-spectra-a.dat};
\end{axis}
\end{tikzpicture}\par
\begin{tikzpicture}
\begin{axis}[dept, xmin=0, xmax=220, xtick={0,50,100,150,200}, title={DEPT-90}]
\addplot[omMeth] table[x=shift, y=d90] {figdata/out/bachelor-2-nmr-spectra-a.dat};
\end{axis}
\end{tikzpicture}
\end{minipage}\hfill
\begin{minipage}[t]{0.49\linewidth}\centering
{\small benzoate d'éthyle}\par
\begin{tikzpicture}
\begin{axis}[dept, xmin=0, xmax=220, xtick={0,50,100,150,200}, title={découplé}]
\addplot[omDef] table[x=shift, y=dec] {figdata/out/bachelor-2-nmr-spectra-b.dat};
\end{axis}
\end{tikzpicture}\par
\begin{tikzpicture}
\begin{axis}[dept, xmin=0, xmax=220, xtick={0,50,100,150,200}, title={DEPT-135}]
\addplot[omThm] table[x=shift, y=d135] {figdata/out/bachelor-2-nmr-spectra-b.dat};
\end{axis}
\end{tikzpicture}\par
\begin{tikzpicture}
\begin{axis}[dept, xmin=0, xmax=220, xtick={0,50,100,150,200}, title={DEPT-90}]
\addplot[omMeth] table[x=shift, y=d90] {figdata/out/bachelor-2-nmr-spectra-b.dat};
\end{axis}
\end{tikzpicture}
\end{minipage}

{\small Spectres du \ce{^{13}C} découplés (déplacements et hauteurs mesurés, données PubChem) avec les
réponses \omterm{def:b2:structure-determination:dept}{DEPT} (signes d'après la \cref{prop:b2:structure-determination:dept}, hauteurs schématiques).
Butan-2-one~: C=O 209.3 (absent en \omterm{def:b2:structure-determination:dept}{DEPT}), \ce{CH2} 36.9 (vers le bas), \ce{CH3} 29.4 et 7.9 (vers le
haut). Benzoate d'éthyle~: le C=O à 166.5 et le carbone du cycle qui porte l'ester, à 130.6,
disparaissent~; trois CH \omterm{def:b2:huckel:aromatic}{aromatiques} (132.8, 129.6, 128.3) restent en DEPT-90~; \ce{OCH2} à 60.9 pointe
vers le bas.}
% figdata: bachelor-2/nmr-spectra.py parts a, b
\end{omfigure}

\section{Combiner les techniques}

Chaque technique répond à sa propre question, et les réponses s'assemblent dans un ordre fixe~: la
formule d'abord, parce qu'elle borne tout le reste~; puis les groupes fonctionnels et le squelette~;
puis les enchaînements.

\begin{method}[Résoudre une inconnue]\label{met:b2:structure-determination:solve}
\begin{enumerate}
  \item Formule~: \omterm{def:b2:mass-spec-atomic:molecular-ion}{ion moléculaire}, $M+1$ pour le nombre de carbones, \omterm{def:b2:mass-spec-atomic:isotope-pattern}{massifs isotopiques}, \omterm{def:b2:mass-spec-atomic:nitrogen-rule}{règle de
        l'azote}~; une masse exacte si elle est disponible. Nombre d'insaturations.
  \item Infrarouge~: \ce{C=O} (et de quelle sorte), \ce{O-H}, \ce{N-H}, \ce{C#N}, \ce{C-H} \omterm{def:b2:huckel:aromatic}{aromatique}
        ou alcénique au-dessus de \qty{3000}{cm^{-1}}.
  \item Carbone~13 et \omterm{def:b2:structure-determination:dept}{DEPT}~: nombre de sortes de carbone (symétrie), leurs classes d'après les
        déplacements, et C, CH, \ce{CH2}, \ce{CH3} d'après la \omterm{def:b2:structure-determination:dept}{DEPT}. Comparer avec la formule.
  \item RMN du proton~: intégrations, déplacements, multiplicités~; voisins d'après la règle des $n + 1$ uplets.
  \item Assembler les fragments~; écrire toutes les structures candidates.
  \item Confronter chaque candidate à chaque donnée, y compris aux fragments du \omterm{def:b2:mass-spec-atomic:spectrum}{spectre de masse}~; garder
        celle qui les explique toutes.
\end{enumerate}
\end{method}

\begin{omfigure}
\begin{tikzpicture}[x=1cm, y=1cm, st/.style={draw=black!70, rounded corners=2pt, fill=white,
  font=\small, align=center, minimum height=0.9cm, inner sep=3pt, text width=2.3cm}]
\node[st, fill=omDef!10] (ms) at (0,0) {SM\\formule, $M+1$};
\node[st] (dbe) at (2.9,0) {nombre\\d'insaturations};
\node[st, fill=omDef!10] (ir) at (5.8,0) {IR\\groupes fonctionnels};
\node[st, fill=omDef!10] (c13) at (8.7,0) {\ce{^{13}C}, DEPT\\squelette};
\node[st, fill=omDef!10] (h1) at (8.7,-1.8) {\ce{^1H}\\enchaînements};
\node[st] (cand) at (5.8,-1.8) {structures\\candidates};
\node[st, fill=omThm!10] (chk) at (2.9,-1.8) {vérifier chaque\\donnée};
\node[st] (ans) at (0,-1.8) {structure};
\draw[->, thick] (ms) -- (dbe); \draw[->, thick] (dbe) -- (ir); \draw[->, thick] (ir) -- (c13);
\draw[->, thick] (c13) -- (h1); \draw[->, thick] (h1) -- (cand); \draw[->, thick] (cand) -- (chk);
\draw[->, thick] (chk) -- (ans);
\draw[->, thick, densely dashed] (chk.south) -- ++(0,-0.5) -| node[pos=0.25, below, font=\footnotesize]
  {une donnée inexpliquée~: retour aux candidates} (cand.south);
\end{tikzpicture}
\par\vspace{4pt}

{\small L'ordre de travail pour une inconnue (\cref{met:b2:structure-determination:solve}). La dernière
étape est celle qu'on omet le plus souvent~: une structure qui laisse un pic inexpliqué n'est pas encore
la réponse.}
\end{omfigure}

\section{Cas traités}

\begin{example}[Une cétone, C$_5$H$_{10}$O]\label{ex:b2:structure-determination:ketone}
Une inconnue a $M = 86$, une bande IR vers \qty{1715}{cm^{-1}}, des raies du carbone à 208.9, 45.7, 29.8,
17.4 et
13.7~ppm, et des signaux du proton à 2.41 (t, 2H), 2.13 (s, 3H), 1.61 (sextuplet, 2H) et 0.91 (t, 3H).
Une insaturation, une cétone saturée (pas de proton aldéhydique vers 9.7). Cinq carbones, cinq raies~:
pas de symétrie. Le singulet de trois protons à 2.13 est un \ce{CH3} voisin du carbonyle~; la suite
triplet, sextuplet, triplet est un groupe propyle dont le \ce{CH2} à 2.41 touche le carbonyle. C'est la
pentan-2-one,
\ce{CH3COCH2CH2CH3}. Son isomère symétrique, la pentan-3-one, présenterait trois raies du carbone et pas
de singulet~; son \omterm{def:b2:mass-spec-atomic:spectrum}{spectre de masse} présente l'ion de McLafferty à 58, absent pour la pentan-3-one (\cref{ch:b2:mass-spec-atomic}).
% ledger: c13:pentan-2-one, nmr:pentan-2-one, ir:CO-ketone, ms:pentan-2-one
\end{example}

\begin{example}[Un ester aromatique, C$_9$H$_{10}$O$_2$]\label{ex:b2:structure-determination:ester}
Une inconnue de formule \ce{C9H10O2} (cinq insaturations) présente une bande de carbonyle et sept raies
du carbone~: 166.5 (C), 132.8 (CH), 130.6 (C), 129.6 (CH), 128.3 (CH), 60.9 (\ce{CH2}) et 14.3
(\ce{CH3}). Son spectre du proton~: 8.05 (m, 2H), 7.35--7.63 (m, 3H), 4.37 (q, 2H), 1.38 (t, 3H). Un
cycle benzénique monosubstitué (quatre raies de carbones \omterm{def:b2:huckel:aromatic}{aromatiques}, dont deux pour deux carbones
chacune~; cinq protons \omterm{def:b2:huckel:aromatic}{aromatiques}) rend compte de quatre insaturations, le carbonyle de la cinquième. La
paire quadruplet--triplet est un groupe éthyle~; son \ce{CH2} à 4.37 (et à 60.9 en carbone) est lié à
un oxygène. Le carbonyle d'ester à 166.5 est lié au cycle~: les deux protons à 8.05, voisins de lui, sont
déplacés vers les champs faibles. C'est le benzoate d'éthyle, \ce{C6H5COOCH2CH3}.
% ledger: c13:ethylbenzoate, nmr:ethylbenzoate
\end{example}

\section{Pièges}

\begin{itemize}
  \item \emph{Protons échangeables} (\ce{O-H}, \ce{N-H})~: ils donnent des signaux larges à des
        déplacements variables, sont souvent non couplés à leurs voisins, et disparaissent quand on agite
        l'échantillon avec une goutte de \ce{D2O}.
  \item \emph{Signaux superposés}~: deux carbones d'un cycle benzénique, ou plusieurs \ce{CH2} d'une chaîne,
        peuvent partager une raie~; compter les raies en regard de la formule avant de conclure sur la symétrie.
  \item \emph{Protons diastéréotopes}~: les deux protons d'un \ce{CH2} voisin d'un stéréocentre ne sont pas
        équivalents, peuvent avoir des déplacements différents et se coupler entre eux.
  \item \emph{Spectres du second ordre}~: quand deux protons couplés ont des déplacements plus proches que
        quelques fois leur constante de couplage, les multiplets penchent l'un vers l'autre et la règle
        des $n+1$ uplets est en défaut~; un spectromètre à champ plus élevé rétablit des motifs simples.
\end{itemize}

\begin{safety}{\ghs[5mm]{GHS02}\,\ghs[5mm]{GHS07}\,\ghs[5mm]{GHS08}}
La butan-2-one est inflammable. L'aimant d'un spectromètre de RMN est toujours sous champ~: pas d'outils
en acier, de bouteilles de gaz ni de cartes magnétiques à proximité, et pas d'accès aux porteurs de
stimulateur cardiaque.
% ledger: ghs:butanone, ghs:benzylacetate
\end{safety}

\section{Exercices}

\begin{exercise}[$\star$]\label{exo:b2:structure-determination:1}
Combien de raies présente le spectre du \ce{^{13}C} découplé de chaque xylène (1,2-, 1,3- et
1,4-diméthylbenzène)~?
\end{exercise}

\begin{exercise}[$\star$]\label{exo:b2:structure-determination:2}
Prévoir les spectres DEPT-135 et DEPT-90 du butan-2-ol.
\end{exercise}

\begin{exercise}[$\star$]\label{exo:b2:structure-determination:3}
À quelle sorte de carbone appartient le plus probablement une raie à 209, 170, 129, 64 et 14~ppm~?
\end{exercise}

\begin{exercise}[$\star$]\label{exo:b2:structure-determination:4}
Calculer le nombre d'insaturations de \ce{C8H8O2} et proposer une structure contenant un cycle benzénique.
\end{exercise}

\begin{exercise}[$\star\star$]\label{exo:b2:structure-determination:5}
Un liquide \ce{C4H8O} présente une bande IR à \qty{1715}{cm^{-1}} et des raies du carbone à 209, 37, 29
et 8~ppm (DEPT-135~: 37 vers le bas, 29 et 8 vers le haut, 209 absente). L'identifier.
\end{exercise}

\begin{exercise}[$\star\star$]\label{exo:b2:structure-determination:6}
Comment distinguer la pentan-2-one de la pentan-3-one par la seule RMN du \ce{^{13}C}~? Par la RMN du
proton~? Par \omterm{def:b2:mass-spec-atomic:spectrum}{spectrométrie de masse}~?
\end{exercise}

\begin{exercise}[$\star\star$]\label{exo:b2:structure-determination:7}
Un ester a $M = 88$ et des signaux du proton à 4.12 (q, 2H), 2.04 (s, 3H) et 1.26 (t, 3H) (données de
l'exercice). L'identifier.
\end{exercise}

\begin{exercise}[$\star\star$]\label{exo:b2:structure-determination:8}
Distinguer le 1,2-dichlorobenzène du 1,4-dichlorobenzène par RMN du \ce{^{13}C}.
\end{exercise}

\begin{exercise}[$\star\star$]\label{exo:b2:structure-determination:9}
Dans le 2-chlorobutane, pourquoi les deux protons du groupe \ce{CH2} ne sont-ils pas équivalents~? Quel
en est l'effet sur le spectre du proton~?
\end{exercise}

\begin{exercise}[$\star\star\star$]\label{exo:b2:structure-determination:10}
Une inconnue \ce{C9H10O} présente une bande IR à \qty{1690}{cm^{-1}}, des raies du carbone à 200.8 (C),
137.0 (C), 132.9 (CH), 128.6 (CH), 128.0 (CH), 31.8 (\ce{CH2}) et 8.3 (\ce{CH3}), et des signaux du
proton à 7.95 (m,
2H), 7.40--7.55 (m, 3H), 2.98 (q, 2H) et 1.22 (t, 3H) (données de l'exercice). L'identifier, et expliquer
le faible nombre d'onde du carbonyle.
\end{exercise}

\begin{exercise}[$\star\star\star$]\label{exo:b2:structure-determination:11}
Quatre isomères \ce{C5H10O2}~: l'acide pentanoïque, le butanoate de méthyle, le propanoate d'éthyle et
l'éthanoate de propyle. Donner pour chacun une particularité de son spectre IR ou de RMN du proton qui le
distingue des trois autres.
\end{exercise}

\begin{exercise}[$\star\star\star$]\label{exo:b2:structure-determination:12}
Un rapport attribue la raie à 166.5~ppm du benzoate d'éthyle au carbone du cycle qui porte l'ester, et
la raie à 130.6~ppm au carbone du carbonyle. Montrer, à l'aide des données \omterm{def:b2:structure-determination:dept}{DEPT} et du diagramme des
déplacements, que cette attribution ne peut pas être juste.
\end{exercise}

\section{Problème~: quatre spectres au parfum de jasmin}

\begin{problem}[{Problème du week-end --- la formule tirée de l'ion moléculaire, la bande infrarouge du
carbonyle, le carbone~13 et la DEPT, le spectre du proton, et une structure confrontée aux fragments}]
\label{pb:b2:structure-determination:1}
Un liquide incolore à l'odeur fleurie de jasmin donne les données suivantes. \omterm{def:b2:mass-spec-atomic:spectrum}{Spectre de masse} (NIST
WebBook)~: $m/z$ 150 (31.7), 151 (3.1), 108 (100), 91 (71.7), 90 (40.6), 43 (37.6). Infrarouge (données
de l'exercice)~: bandes intenses à 1740 et \qty{1230}{cm^{-1}}, bandes faibles juste au-dessus de
\qty{3000}{cm^{-1}}, pas de
bande large vers \qty{3300}{cm^{-1}}. Carbone~13, CDCl$_3$ (données PubChem)~: 170.70, 136.14, 128.56,
128.24, 66.24, 20.82~ppm. Proton, CDCl$_3$, \qty{90}{MHz}~: 7.33 (m, 5H), 5.09 (s, 2H), 2.06 (s, 3H).
% ledger: use:benzylacetate, ms:benzylacetate, c13:benzylacetate, nmr:benzylacetate, ir:CO-ester, iso:C12, iso:C13, mass:C12, mass:H1, mass:O16

\begin{center}
\begin{tikzpicture}
\pgfplotsset{dept/.style={width=0.8\linewidth, height=3.2cm, ycomb, x dir=reverse, ymin=0, ymax=1.15,
  xmin=0, xmax=200, ytick=\empty, axis y line=none, axis x line=bottom, x axis line style={-}, tick label style={font=\tiny},
  title style={font=\scriptsize, at={(0.02,0.82)}, anchor=west}, every axis plot/.append style={thick}}}
\begin{axis}[dept, title={découplé}, xtick={0,20,40,60,80,100,120,140,160,180,200}]
\addplot[omDef] table[x=shift, y=dec] {figdata/out/bachelor-2-nmr-spectra-c.dat};
\end{axis}
\end{tikzpicture}
\end{center}
% figdata: bachelor-2/nmr-spectra.py part c

\medskip
\noindent\textbf{Partie I --- La formule.}
\begin{enumerate}
  \item À partir du rapport 151/150, estimer le nombre de carbones.
  \item Que dit la \omterm{def:b2:mass-spec-atomic:nitrogen-rule}{règle de l'azote}~?
  \item Avec neuf carbones, que reste-t-il de la masse 150~? Proposer une formule comportant de l'oxygène.
  \item Pourquoi \ce{C10H14O}, lui aussi de masse nominale 150, est-il moins probable~?
  \item Calculer le nombre d'insaturations.
  \item Calculer la \omterm{def:b2:mass-spec-atomic:exact-mass}{masse monoisotopique} de la formule retenue.
\end{enumerate}

\medskip
\noindent\textbf{Partie II --- Infrarouge.}
\begin{enumerate}[resume]
  \item Attribuer la bande à \qty{1740}{cm^{-1}}. Quelle famille sa position suggère-t-elle~?
  \item Attribuer la bande à \qty{1230}{cm^{-1}}.
  \item Qu'indiquent les bandes faibles au-dessus de \qty{3000}{cm^{-1}}~?
  \item Qu'exclut l'absence de bande large vers \qty{3300}{cm^{-1}}~?
  \item Le carbonyle est-il conjugué avec le cycle~? Expliquer d'après son nombre d'onde.
\end{enumerate}

\medskip
\noindent\textbf{Partie III --- Carbone~13 et \omterm{def:b2:structure-determination:dept}{DEPT}.}
\begin{enumerate}[resume]
  \item Attribuer chaque raie à une sorte de carbone.
  \item Quelle raie pointe vers le bas en DEPT-135~?
  \item Quelles raies disparaissent en DEPT-135~?
  \item Quelles raies subsistent en DEPT-90~?
  \item Combien de sortes de carbone un cycle benzénique monosubstitué possède-t-il~? Expliquer pourquoi
        seules deux raies de CH \omterm{def:b2:huckel:aromatic}{aromatiques} apparaissent.
  \item La raie à 128.24 est la plus haute. Représente-t-elle le plus de carbones~?
\end{enumerate}

\medskip
\noindent\textbf{Partie IV --- RMN du proton et structure.}
\begin{enumerate}[resume]
  \item Attribuer les trois signaux du proton.
  \item Pourquoi les signaux à 5.09 et 2.06 sont-ils des singulets~?
  \item Pourquoi le \ce{CH2} à 5.09 est-il si déplacé vers les champs faibles~?
  \item Assembler la structure et la nommer.
  \item Son isomère, le 2-phényléthanoate de méthyle, \ce{C6H5CH2COOCH3}, a la même formule. Quelle donnée
        l'exclut~?
  \item Expliquer les fragments à 108, 91 et 43.
  \item Combien de raies présenterait un spectre du \ce{^{13}C} découplé parfaitement résolu du composé,
        et combien d'entre elles pointent vers le bas en DEPT-135~?
\end{enumerate}
\end{problem}
